(1) 理想Michelson干涉仪以波长 $\lambda_0$ 的单色光工作，两臂几何长度为 $L_1,L_2$，光在空气中传播，求往返光程差。
(2) 推导两回波重叠输出的干涉强度关于臂长差的表达式。
(3) 将第一臂镜面平移微小距离 $\Delta x$，求引起的相位变化。
(4) 若探测器统计明暗条纹经过数 $N$，反推镜面位移，并给出移动方向不能从单一强度计数判断的原因。
